\subsection{Particle Physics}
\newpage
\begin{itemize}
	\item Transverse momentum \pT:\@\ momentum transverse to the beam axis.
	\item Transverse Energy \ET:\@\ Energy transverse to beam axis, deposited in calorimeter.
	\item Scalar Fields
	\item Classical Scattering
	\item ATLAS Data
	\begin{itemize}
		\item Run 1: 2010--2012
		\item Run 2: 2015--2018
		\item Run 3: 2022--2026
	\end{itemize}
	\item Cutflow: Volume of phase space (number of events) after each cut is applied. Used to visualize the cut restrictions as they are applied. 
	\item Validation Region / Control Region: A selection of data space that very close to signal selection region. Used to `validate' the background region. Usually just has same cuts as signal region with exception of one cut inverted.
	\item Quantization
	\item On-shell: 
	\item Branching Ratio: The fraction of particles that decay via a particular decay mode.
	\item Cross Section: A measure of the probability that a specific process will occur when particles
	      collide. It is typically measured in units of area, such as barns (1 barn = $\rm 10^{-28} m^2$ = $\rm 10^{-24} cm^2$).
	\item Luminosity: A measure of the number of particles per unit area per unit time in a collider. It is used to calculate the event rate for a given process.  \(L=\oo{\sigma}\der{N}{t}\). 
	\item Integrated Luminosity: The total luminosity accumulated over a specific period of time. It is often expressed in inverse femtobarns (\(\rm fb^{-1}\)). This is the quantity that relates to the amount of data we have. 
	\item Event Topologies.
\end{itemize}
\subsection{Width}
When particles decay immeasurably fast (\W, \Z, \h, t), width $\Gamma$ is used instead of mean lifetime $\tau$. 
\subsection{Higgs}
The Higgs boson \h\ explains how the vector bosons get their mass (why it's \W, \Z and $\gamma$) instead of all or none having mass. 
The \h\ also explains how fermions get their mass, and there is a correspoondance between Newtons first law. Interacting with the higgs field makes it difficult for a body to accelerate, which is analogous to inertia. do I just need more text at the end?
\subsection{ZZ}
Z mass gets smeared in reconstruction, rapidity \emph{does not}. I'm not sure why, but this is useful for truth matching with reconstructed events. 

%\subsection{Feynman Rules}
%	\tikzfeynmanset{
%		every vertex={red},
%		every particle={blue},
%		every blob={draw=black, pattern color=black},
%		}
%	\begin{align*}
%		%Incoming anti/fermions
%		\feynmandiagram[small, baseline=(v1.base), horizontal=v1 to v2]
%		{
%			v1 [color=black, particle=$f$] -- [fermion] v2[blob]-- [draw=none] v3 [particle={\phantom{$f$}}, color=white]
%		};
%		&\quad\longleftrightarrow\quad
%		u(p,s)
%		\\
%		\feynmandiagram[small, baseline=(v1.base), horizontal=v1 to v2]
%		{
%			v1 [color=black, particle=$f$] -- [anti fermion] v2[blob]-- [draw=none] v3 [particle={\phantom{$f$}}, color=white]
%		};
%		&\quad\longleftrightarrow\quad
%		\overline{v}(p,s)
%		\\
%		% Outgoing anti/fermions
%		\feynmandiagram[small, baseline=(v1.base), horizontal=v1 to v2]
%		{
%			v3 [draw=none] -- [draw=none] v1 [blob] -- [fermion] v2 [color=black, particle=$f$]
%		};
%		&\quad\longleftrightarrow\quad
%		\overline{u}(p,s)
%		\\
%		\feynmandiagram[small, baseline=(v1.base), horizontal=v1 to v2]
%		{
%			v3 [draw=none] -- [draw=none] v1 [blob] -- [anti fermion] v2 [color=black, particle=$f$]
%		};
%		&\quad\longleftrightarrow\quad
%		v(p,s)
%		\\
%		% Photon fermion interaction
%		\feynmandiagram[small, baseline=(v1.base), horizontal=v1 to v2]{
%		v1 [particle=$\mu$, color=black] -- [photon] v2,
%		[particle=$f$] a -- [fermion] v2 -- [fermion] b [particle=$f$],
%		};
%		&\quad\longleftrightarrow\quad
%		-iQ_fe\gamma^{\mu}
%	\end{align*}%